x\chapter{Introduktion}
\label{chap:introduction}

\section{Baggrund}
\label{sec:baggrund}

Arkyn Studios er en dansk softwarevirksomhed, der udvikler digitale løsninger til enterprise-virksomheder med fokus på vedligeholdelse af maskiner og udstyr i produktionsmiljøer. Virksomhedens kerneprodukt er en mobil applikation, der understøtter teknikernes daglige arbejde på fabriksgulvet — fra registrering af opgaver til dokumentation af udført vedligehold.

I takt med at produktionsmiljøer automatiseres og specialiseres, oplever mange virksomheder et stigende behov for digital støtte til de medarbejdere, der udfører vedligeholdelsesopgaverne i praksis. Arkyn Studios' kunder har specifikt givet udtryk for, at fremtidens vedligeholdelsesteknikere i stigende grad vil have en mere generalistisk baggrund, og dermed have mindre dybdegående kendskab til den specifikke maskinpark de arbejder med. Dette stiller nye krav til de digitale værktøjer, der skal støtte dem i arbejdet.

Dette projekt er gennemført i tilknytning til et praktikophold hos Arkyn Studios og tager udgangspunkt i en konkret problemstilling identificeret i samarbejde med virksomheden.

\section{Problemformulering}
\label{sec:problemformulering}

En konkret og gentagen udfordring i vedligeholdelsesarbejdet er, at teknikeren hurtigt skal kunne identificere en maskine og finde relevant information om dens vedligeholdelse — særligt i situationer, hvor maskinen ikke er umiddelbart genkendelig, eller hvor teknikeren ikke har erfaring med netop denne maskintype.

I dag sker denne identifikation primært manuelt: teknikeren aflæser maskinens typeskilt eller mærkat og søger derefter i dokumentation eller kontakter en mere erfaren kollega. Denne proces er tidskrævende, fejlbehæftet og afhænger i høj grad af den individuelle teknikers erfaring og viden.

Selvom teknologier som Optical Character Recognition (OCR) og kunstig intelligens eksisterer som separate løsninger, er integrationen af billedbaseret input og AI-assisteret vejledning i en vedligeholdelseskontekst fortsat et åbent og relevant problemfelt.

\begin{infobox}[Problemformulering]
Hvordan kan en mobilapplikation ved hjælp af OCR og AI-genereret vejledning understøtte vedligeholdelsesteknikere i hurtig identifikation af maskiner og adgang til relevant vedligeholdelsesinformation?
\end{infobox}

\section{Projektmål}
\label{sec:projektmaal}

På baggrund af problemformuleringen er følgende konkrete mål fastsat for projektet:

\begin{enumerate}
  \item Udvikle en prototype på en mobil feature til iOS, der muliggør billedbaseret identifikation af maskiner ved optagelse af typeskilte og maskinlabels
  \item Implementere en pipeline der via OCR udtrækker struktureret information fra billeder og korte videoer
  \item Integrere en AI-baseret service, der på baggrund af den udtrukne information genererer forståelig og handlingsorienteret vedligeholdelsesvejledning
  \item Designe et simpelt og effektivt brugerflow baseret på etablerede usability-principper, der fungerer i en industriel arbejdskontekst
  \item Evaluere systemets tekniske funktionalitet og brugbarhed i samarbejde med Arkyn Studios
\end{enumerate}

\section{Afgrænsning}
\label{sec:afgrænsning}

Projektet er afgrænset til udviklingen af en prototype, og følgende er bevidst udeladt af scope:

\textbf{Ingen ERP/SAP-integration:} Systemet henter ikke information fra virksomhedsspecifikke systemer. Al vejledning genereres på baggrund af generel industriel viden fra AI-modellen.

\textbf{iOS only:} Applikationen er udelukkende udviklet til iOS (iPhone og iPad) og understøtter ikke Android i den nuværende prototype.

\textbf{Ingen autentificering:} Prototypen implementerer ikke brugerlogin eller rollestyring, da fokus er på kerneflowet frem for sikkerhedsinfrastruktur.

\textbf{Begrænset maskindomæne:} Systemet er ikke specialiseret til én bestemt maskintype, men demonstreres med generelle industrielle maskiner og udstyr.

